\begingroup
\setlength{\tabcolsep}{4pt}
\renewcommand{\arraystretch}{1.13}
\renewcommand{\arraystretch}{1.13}
\small
\begin{longtable}{lrr}
\caption{Characteristics of the harmonized multiple-myeloma analysis cohorts.}\label{tab:table5}\\
\toprule
Characteristic & EloKRd & UCMM \\ \midrule
\endfirsthead
\multicolumn{3}{l}{\tablename\ \thetable\ (continued)}\\
\toprule
Characteristic & EloKRd & UCMM \\ \midrule
\endhead
\midrule\multicolumn{3}{r}{Continued on next page}\\\endfoot
\bottomrule\endlastfoot
Participants & 30 & 253 \\
Age (years), median (IQR) & 62.8 (57.4, 70.2) & 61.3 (54.4, 67.6) \\
Male & 21 (70.0) & 130 (51.4) \\
Race &  &  \\
\quad White & 23 (76.7) & 164 (64.8) \\
\quad Black & 4 (13.3) & 76 (30.0) \\
\quad Other & 3 (10.0) & 13 (5.1) \\
Hispanic or Latino & 1 (3.3) & 13 (5.1) \\
High-risk cytogenetics & 15 (50.0) & 58 (22.9) \\
ASCT & 9 (30.0) & 203 (80.2) \\
Progression or death & 13 (43.3) & 159 (62.8) \\
All-cause death & 10 (33.3) & 60 (23.7) \\
\end{longtable}
\noindent\begin{minipage}{\linewidth}\footnotesize Entries are $n$ (\%) except where indicated. IQR, interquartile range; ASCT, autologous stem-cell transplantation. ASCT is a postinduction characteristic and is not presented as a pretreatment baseline variable. All analysis covariates were complete in both cohorts.\end{minipage}
\endgroup
